% Architecture documentation for todo-backend.
% Build with docs/build.sh, which renders the diagrams to build/figures/*.pdf first.
\documentclass[a4paper,11pt]{article}

\usepackage[T1]{fontenc}
\usepackage{lmodern}
\usepackage{microtype}
\usepackage[margin=2cm]{geometry}
\usepackage{graphicx}
\usepackage{pdflscape}
\usepackage{booktabs}
\usepackage{parskip}
\usepackage[font=small,labelfont=bf,hypcap=false]{caption}
\usepackage[hidelinks]{hyperref}

\graphicspath{{build/figures/}}

% build.sh writes the git revision here; fall back gracefully when run by hand.
\IfFileExists{build/version.tex}{\input{build/version.tex}}{\newcommand{\docversion}{unknown}}

\hypersetup{
  pdftitle={todo-backend architecture},
  pdfsubject={Architecture diagrams},
}

% Page dimensions in portrait orientation, saved before any landscape page swaps them.
\newlength{\portraitwidth}
\newlength{\portraitheight}
\AtBeginDocument{%
  \setlength{\portraitwidth}{\textwidth}%
  \setlength{\portraitheight}{\textheight}%
}

% Room left for a section heading, a short paragraph and the caption.
\newcommand{\textallowance}{9\baselineskip}

% \diagram{file}{caption}: a figure scaled to fill the rest of a portrait page.
% Not a float, so it stays on the same page as its section text.
\newcommand{\diagram}[2]{%
  \begin{center}
    \includegraphics[width=\portraitwidth,
                     height=\dimexpr\portraitheight-\textallowance\relax,
                     keepaspectratio]{#1}
    \captionof{figure}{#2}
  \end{center}%
}

% \widediagram{file}{caption}: the same for use inside a landscape environment.
\newcommand{\widediagram}[2]{%
  \begin{center}
    \includegraphics[width=\portraitheight,
                     height=\dimexpr\portraitwidth-\textallowance\relax,
                     keepaspectratio]{#1}
    \captionof{figure}{#2}
  \end{center}%
}

\title{\texttt{todo-backend}\\[0.3em]\large Architecture documentation}
\author{}
\date{\today\\[0.3em]\small Revision \texttt{\docversion}}

\begin{document}

\maketitle
\thispagestyle{empty}

\section*{Overview}

\texttt{todo-backend} is a small JSON API for a todo app: users register, log in, and manage
their own todo entries. It is written in Rust with \textbf{axum} and \textbf{tokio}, stores
data in \textbf{PostgreSQL} through \textbf{sqlx}, hashes passwords with \textbf{argon2id}, and
authenticates requests with \textbf{HS256 JWT} bearer tokens. It runs as two containers, the
app and its database, via Docker Compose.

The diagrams in this document are generated from the PlantUML sources in
\texttt{docs/architecture/}; edit those and rebuild with \texttt{docs/build.sh}.

\begin{center}
\small
\begin{tabular}{@{}rlll@{}}
  \toprule
  \# & Diagram & Type & Shows \\
  \midrule
  1  & System context   & use case   & Clients, the API and PostgreSQL \\
  2  & Deployment       & deployment & Compose services, network, volume, test DB \\
  3  & Components       & component  & Crate modules and their dependencies \\
  4  & Request pipeline & activity   & Tower layers, extractors, error mapping \\
  5  & Register         & sequence   & \texttt{POST /auth/register} \\
  6  & Login            & sequence   & \texttt{POST /auth/login} \\
  7  & Update a todo    & sequence   & \texttt{PATCH /todos/\{id\}} \\
  8  & Data model       & ER         & \texttt{users} and \texttt{todos} tables \\
  9  & Models           & class      & Request and response types, validation \\
  10 & Errors and state & class      & \texttt{AppError}, \texttt{Config}, \texttt{AppState} \\
  11 & Lifecycle        & activity   & Startup, migrations, graceful shutdown \\
  12 & Test layout      & component  & Test suites and what each exercises \\
  \bottomrule
\end{tabular}
\end{center}

\clearpage
\tableofcontents
\listoffigures
\clearpage

\section{System context}
The API serves two kinds of client: scripts and tools calling it directly, and browser
applications on origins listed in \texttt{CORS\_ALLOWED\_ORIGINS}. All persistent state lives in
PostgreSQL, reached through a pool of at most ten connections.
\diagram{system-context}{System context}

\begin{landscape}
\section{Deployment}
Docker Compose runs the app and its database on a private network; only the app's port is
published. The \texttt{db-test} service, enabled with the \texttt{test} profile, is a throwaway
database on \texttt{tmpfs} for \texttt{cargo test} and sqlx's compile-time query checks.
\widediagram{deployment}{Deployment (\texttt{docker-compose.yml})}
\end{landscape}

\begin{landscape}
\section{Components}
\texttt{main.rs} loads configuration, initialises tracing, builds the shared
\texttt{AppState} and hands it to \texttt{build\_app}. Routes depend on extractors, request
models and repositories; only the repositories talk to the database. \texttt{AppError} is the
single error type shared by every layer.
\widediagram{components}{Crate modules and dependencies}
\end{landscape}

\section{Request pipeline}
Every request passes through tracing, CORS and the 64\,KiB body limit before routing. Handler
arguments are custom extractors that turn malformed input into an \texttt{AppError}, so every
error response has the same \texttt{\{"error": "..."\}} shape.
\diagram{request-pipeline}{Request pipeline}

\begin{landscape}
\section{Registration}
Passwords are hashed with argon2id on tokio's blocking thread pool so the CPU-heavy work does
not stall the async runtime. A duplicate username surfaces as a unique-constraint violation and
is reported as \texttt{409 Conflict}.
\widediagram{seq-register}{\texttt{POST /auth/register}}
\end{landscape}

\begin{landscape}
\section{Login}
When the username is unknown, the password is still verified against a dummy hash, so the
response time does not reveal which usernames exist. Successful logins receive an HS256 token
whose lifetime is \texttt{JWT\_TTL\_SECONDS}.
\widediagram{seq-login}{\texttt{POST /auth/login}}
\end{landscape}

\begin{landscape}
\section{Authenticated todo request}
\texttt{PATCH /todos/\{id\}} is representative of all todo routes: the \texttt{AuthUser}
extractor verifies the token, the request body is validated into a change set, and the query
filters on both the todo id and the caller's id. Another user's todo is therefore
indistinguishable from a missing one.
\widediagram{seq-todo-update}{\texttt{PATCH /todos/\{id\}}}
\end{landscape}

\begin{landscape}
\section{Data model}
Two tables, created by the migrations in \texttt{migrations/} and applied automatically on
startup. Deleting a user cascades to their todos.
\widediagram{data-model}{Database schema}
\end{landscape}

\section{Request and response models}
Request types are validated into separate, already-checked types before reaching the
repositories, so a repository never sees unvalidated input. \texttt{User} carries the password
hash and is never serialised; responses use \texttt{UserResponse}.
\diagram{models}{Request and response models}

\begin{landscape}
\section{Errors, configuration and shared state}
\texttt{AppError} maps each failure to an HTTP status; server-side details are logged rather
than returned. \texttt{AppState} is cloned into every handler and holds the connection pool and
the configuration loaded at startup.
\widediagram{errors-and-state}{Errors, configuration and shared state}
\end{landscape}

\section{Process lifecycle}
Configuration errors abort startup with a message. On SIGTERM (as sent by \texttt{docker stop})
or Ctrl-C the server stops accepting connections, lets in-flight requests finish, and closes
the pool.
\diagram{lifecycle}{Process lifecycle}

\clearpage
\section{Test layout}
Unit tests cover pure logic. Repository and API tests run against a fresh database per test via
\texttt{\#[sqlx::test]}, with the API exercised in-process through \texttt{tower}. End-to-end
tests drive the real Docker Compose stack over HTTP and run in CI.
\diagram{test-layout}{Test suites}

\end{document}
